% Full primitive-region reduction, nested substitution and final D-image construction.

\section{Completion of Kamp's expressive-completeness construction}
\label{sec:annex-p09-kamp-iii2-completion}

\begin{annexitem}{annex:p9.kamp.one.sided.closure}{thm:p09-kamp-one-sided-closure}{Both guarded closures with their actual elementary histories}
\begin{proof}
Suppose \(\rho(x,w)\) has a supplied representation
\[
 \mathsf{Com}\models x<w\to
                  (\rho(x,w)\longleftrightarrow\bigvee_i\theta_i(x,w)),
\]
where each \(\theta_i\) belongs to \(\operatorname{Ep}(x,w)\).
The existential guard supplies exactly the hypothesis under which the
representation is valid. Distributing the existential over this finite
disjunction gives
\[
 \exists x\,(x<w\land\rho(x,w))
       \longleftrightarrow\bigvee_i\exists x\,(x<w\land\theta_i(x,w))
 \tag{1}
\]
over \(\mathsf{Com}\). Each disjunct on the right has an
actual \(\operatorname{E}(w)\) presentation by the elementary
past-closure construction of
\cref{thm:p09-kamp-elementary-intervals}. That construction starts
with the guarded basic past seed and applies the exact finite
substitution history of its interval word. The outer guard is unchanged
and the endpoint variable is bound as intended. Taking the finite
disjunction of these histories gives an elementary history for (1).
If the given list is empty, use \(F^K_w\).

For a supplied \(I(w,x)\) representation, obtain the elementary
labelled word of each past disjunct. Construct its elementary future
presentation with the same ordered word. On every ordered pair, both
presentations have exactly the same finite realizing sequences, by
the elementary-word theorem. Thus this replacement is valid under
the actual future guard \(w<x\). Distribute the guarded
existential over the resulting future list and use the elementary
future-closure construction for each member of \(\operatorname{Ef}(w,x)\).
Its basic skeleton is a guarded future seed and its supplied
substitution history produces a genuine member of \(\operatorname{E}(w)\).
The disjunction of all outputs proves the second assertion.

Neither argument assumes that an elementary past formula is literally
an elementary future formula or that a future closure of its past
spelling is an elementary seed by definition. The same-word conversion
and the actual closure histories supply the required syntax.
Missing past or future points make the corresponding existential
false, exactly as the constructed formula does. No new witness is
created to satisfy a guard \cite[ch.~III, pp.~67--68, 93--94]{Kamp1968}.
\end{proof}
\end{annexitem}

\begin{annexitem}{annex:p9.kamp.primitive.regions}{thm:p09-kamp-primitive-regions}{Every primitive existential and universal case}
\begin{proof}
In an existential past envelope, the matrix is a finite conjunction
of Bet components at \((x,w)\) and At conditions at those
endpoints. Lemma~10 gives a finite elementary interval representation
for that entire conjunction. Applying
\cref{thm:p09-kamp-one-sided-closure} produces its elementary
past existential presentation. In the future case use the ordered
pair \((w,x)\) and the future half of that closure theorem.
Both operations retain all matrix conjuncts and every disjunct of
their interval representation. An empty matrix uses the truth
label, leaving strict existence of a point on the indicated side.

A primitive Bef member with an existential root has exactly this
past form. Indeed its source constituents at the quantified point
are At, earlier Bef or Bet. An earlier Bef constituent would be
a proper regional occurrence, excluded by membership in
\(\mathcal D\). Therefore only At and Bet remain. The
same reasoning gives the future envelope for an existential Aft
root. Retention in the staged grammar changes no formula spelling
or outermost constructor.

For a primitive universal past member, write its source matrix as
\(\Psi(x,w)\), a finite disjunction of At conditions at
\(x\) and Bet conditions on \((x,w)\). Its complement is
\[
 \neg\forall x\,(x<w\to\Psi(x,w))
          \longleftrightarrow\exists x\,(x<w\land\neg\Psi(x,w)).
 \tag{1}
\]
The interval Boolean normalization theorem gives a positive Boolean
expression equivalent to \(\neg\Psi\) over \(\mathsf{Lin}\),
with leaves only in \(\operatorname{Bet}(x,w)\),
\(\operatorname{At}(x)\), or \(\operatorname{At}(w)\).
Put this finite positive expression in disjunctive normal form and
retain every term. Distribute the existential in (1) over all terms.
Each resulting disjunct is an existential envelope of the exact
kind already treated, including any endpoint At factors generated
by normalization. Construct its elementary presentation and disjoin
the finite list. This gives an elementary formula for the complement
of the original universal. Negating that elementary formula is an
allowed \(\operatorname{E}\) constructor and gives the original
universal's presentation over \(\mathsf{Com}\).

For a primitive universal Aft root, negate the guarded future
universal in the same way, with guard \(w<x\). Normalize
its complemented interval matrix on the ordered pair \((w,x)\),
distribute over its finite DNF terms, and apply the future-envelope
construction to each term. Negate the final elementary disjunction.
Every step has the same explicit direction of equivalence as (1)
and uses the same fixed endpoint; no interchange of a witness
operand and an interval operand occurs.

The empty disjunction in a universal is \(F^K\), so its
complement produces the true existential matrix and detects whether
any point exists on that side. Conversely an empty conjunction in
an existential is \(T^K\) and leaves that same strict witness
condition. If Boolean normalization returns an empty DNF list,
its existential is false and its elementary output is \(F^K_w\).
All these cases respect the source's literal constant convention.

The source moves negations to unary literals and distributes finite
Boolean combinations on p.~62. The interval normalization theorem
supplies its full quantifier and endpoint-factor analysis. Its finite
disjunction is retained here; it is not identified with a single
primitive regional formula merely by changing its name. Closure of
\(\operatorname{E}\) under disjunction and negation is precisely
what permits the conclusion. Each construction uses a finite source
history and the already completed Lemmas~8--10, so it terminates
\cite[ch.~III, pp.~61--62, 93--94]{Kamp1968}.
\end{proof}
\end{annexitem}

\begin{annexitem}{annex:p9.kamp.regional.elementary}{thm:p09-kamp-regional-elementary}{Maximal proper occurrences and the complete-order substitution argument}
\begin{proof}
Induct simultaneously on the size of the original finite syntax tree
for members of Bef and Aft with any sole free variable. A member
with no proper occurrence in those classes is primitive and has
the presentation just proved. For a nonprimitive member
\(\varphi(w)\), list the maximal proper occurrences of
that kind. Their syntax subtrees are disjoint: if two overlap
in a syntax tree, one contains the other, contrary to maximality
unless they are the same occurrence. Each has strictly smaller
syntax size than \(\varphi\), so the induction hypothesis
supplies an actual elementary presentation
\(\xi_i(y_i)\) for \(\varphi_i(y_i)\).

After these finitely many recursive outputs have been constructed,
choose distinct letters \(Q_i\) absent from the whole original
formula and from every \(\xi_i\). This is possible because
the combined set of used letters is finite. Replace the selected
occurrences by the atoms \(Q_i(y_i)\). The source grammar is
preserved: a proper regional constituent is replaced by an allowed
At constituent at its sole free argument. All smaller regional
occurrences inside a selected subtree disappear with it, and none
are introduced by a unary atom. Thus the resulting skeleton
\(\varphi'(w)\) has no proper Bef or Aft occurrence and
belongs to \(\mathcal D(w)\). The primitive-region theorem
constructs its elementary presentation \(\xi'(w)\).

We verify the restoration at its exact order scope. Fix any complete
linear structure \(\mathcal M\). For each \(i\), interpret
\(Q_i\) by the set of points satisfying \(\varphi_i\)
at its one free argument. Because the induction equivalence holds
at every point in this complete structure, this is also the truth
set of \(\xi_i\). Each component has exactly one free
variable, so the set does not depend on any unassigned extra
parameter. These new interpretations leave the complete order
unchanged; the uniform equivalence between \(\varphi'\)
and \(\xi'\) therefore applies to the reinterpreted structure.

The proper-substitution identity transfers that equivalence back
to \(\mathcal M\). Substituting the original \(\varphi_i\)
into their designated fresh atoms reconstructs \(\varphi\),
up to alphabetic variation of internal bound variables. Substituting
\(\xi_i\) yields the constructed elementary formula. These
two restorations agree because the inserted formulas have the
same truth set at every argument in \(\mathcal M\).
This proves the required \(\mathsf{Com}\) equivalence;
no \(\mathsf{Lin}\) equivalence of the recursive regional
components has been assumed.

To exhibit the elementary history, rename the sole free parameter
of each \(\xi_i\) to a common name \(w\), with its
bound names made fresh, and apply the finite sequence of proper
substitutions to \(\xi'(w)\). These are allowed operations
in \(\operatorname{E}(w)\). At an original atom under a
quantifier binding its argument variable \(y_i\), that same
quantifier binds the free argument of the inserted copy. Freshness
protects only its internal bound variables and does not remove
this intended binding. No replacement contains any \(Q_j\)
by the choice of the targets, so later substitutions cannot change
an earlier replacement. Distinct identical occurrences can be
restored through different fresh letters to the same formula,
without losing either occurrence.

Every recursive call concerns a proper original subtree. The
primitive conversion and subsequent substitution histories are finite
and do not recursively reprocess their newly generated outputs.
Thus the procedure terminates. It treats both Bef and Aft in
the simultaneous induction, and the primitive construction supplied
both possible outer quantifiers for both directions. This completes
the source's reduction through maximal proper regional subformulas
\cite[ch.~III, pp.~60--62]{Kamp1968}.
\end{proof}
\end{annexitem}

\begin{annexitem}{annex:p9.kamp.iii2}{thm:p09-kamp-iii2}{The full standard-form, elementary-history and D-image construction}
\begin{proof}
Let \(\varphi\) have exactly the free variable \(w\).
Apply the proved Theorem~III.1 to obtain its literal, unstarred
source-standard formula \(N_\varphi\) on \(\{w\}\).
The equivalence holds throughout \(\mathsf{Lin}\), hence
also throughout \(\mathsf{Com}\). With one free variable
there is one essential variable and no interval between distinct
essential variables. The complete prefix is the coincident
one-variable prefix \(w\doteq w\), true on every linear
order. Every remaining component belongs to
\(\operatorname{At}(w)\), \(\operatorname{Bef}(w)\),
or \(\operatorname{Aft}(w)\). These are exactly the
one-variable component classes in the source's initial reduction
on p.~60.

Replace the true prefix by \(T^K_w\). An At component is
already an elementary seed. For each Bef or Aft component, use
\cref{thm:p09-kamp-regional-elementary} to construct its elementary
presentation over \(\mathsf{Com}\). Retain every conjunction
and every disjunction of the standard form. Each local replacement
is uniformly equivalent at every current point and every unary
interpretation on every complete order, so their finite Boolean
combination preserves the original equivalence. An empty clause
uses the declared truth seed; false cases retain their false
seeds. The output \(\Xi(w)\) has an actual finite
\(\operatorname{E}(w)\) construction history and satisfies
\[
                         \mathsf{Com}\models\varphi(w)\longleftrightarrow\Xi(w).
 \tag{1}
\]

Apply the full elementary-to-temporal construction of
\cref{thm:p09-kamp-iii2-lemma3} to that history. At an At
seed it translates unary letters to their corresponding proposition
letters. At a guarded basic or limiting seed it uses the full
past or future construction of Lemma~1 or Lemma~2. At Boolean
nodes it uses the corresponding Boolean constructors, expanding
disjunction into negation and conjunction. At a proper predicate
substitution node it uses the already proved temporal substitution
operation. Thus it returns a finite \(\eta\in\mathsf{TL}_1\)
with
\[
                         \mathsf{Lin}\models\Xi(w)\longleftrightarrow D(\eta)[w/t_0].
 \tag{2}
\]
Bound variables in the final first-order replacement are made fresh
from \(w\), so only the free current-point variable is replaced.
Restrict (2) to complete orders and combine it with (1). This gives
the displayed theorem for every current point, every unary
interpretation and every complete linear order.

When \(w=t_0\), the output \(D(\eta)\) is a literal
member of the image class \(\mathcal C\), proving the source's
exact Theorem~III.2 statement. The intermediate formula \(\Xi\)
was not assumed to be literally in that image, and no inverse
translation of an arbitrary equivalent first-order spelling is
required. The constructed temporal formula supplies its image member
directly. The first-order semantics of the strict \(S,U\)
translation were already proved in the source's translation step;
the witness remains the first operand and the open interval the
second. Therefore the temporal formula expresses the same uniformly
defined operation on the unary truth sets. This yields the claimed
basis of negation, conjunction, SINCE and UNTIL over complete orders.

For the stated positive finite signature, the construction introduces
no surviving unary letters outside the input signature and the
source's fixed \(P_0\). To check this, the literal normal-form
construction merely repeats input literals and inserts its particular
truth/falsehood formulas. The interval constructions cut or reverse
input words, combine labels, and use elementary limiting histories
whose letters come from those words. Any fresh occurrence letters
used to present elementary words are all replaced by their original
labels. The fresh atom introduced in a smaller-family induction call
is properly replaced by the specified elementary label afterwards.
All such substitutions replace every original occurrence of their
fresh target, and that target occurs in none of their replacement
formulas. The new regional skeleton targets are likewise chosen
outside all replacements and are all restored. Thus no auxiliary
target letter survives at the root.

The same argument can be read as induction on the finite construction
histories. Only the source constants can require an inessential
\(P_0\) not already used by the input. For \(n\geq1\),
that letter is still in the allowed signature
\(P_0,\ldots,P_{n-1}\). The elementary-to-temporal translation
introduces only the corresponding proposition letters, so its output
has the asserted finite-arity support. Equivalence is for every
valuation, so any inessential constant-encoding letter does not
change the represented operation.

Termination follows from the proved constructions: the first-order
normal form processes a finite original syntax; Bet and interval
normalization use finite construction histories; Lemma~8's recursive
families strictly decrease their proved well-founded comparison;
regional restoration recurses on proper original subtrees; and
the final elementary translation is structural recursion on a finite
history. Each nonrecursive stage performs finitely many cuts,
distributions, fusions or substitutions. These procedures construct
one uniform formula rather than choosing a formula after inspecting
a particular model. This establishes the effective construction
claimed here without a size or time estimate.

For the worked universal example, classical quantifier negation and
De Morgan's law give
\(\neg\varphi(w)\longleftrightarrow\exists z\,(z<w\land(Q(z)\lor R(z)))\).
Distributing this existential gives exactly the two displayed guarded
past disjuncts. The formula \(S(q,\tau)\) means that some
strictly earlier point satisfies \(Q\), since \(\tau\)
is true throughout every intervening interval. The formula
\(S(r,\tau)\) has the corresponding meaning for \(R\).
Negating both and conjoining therefore expresses \(\varphi\).
At a first point neither strict witness exists, giving true
universal vacuity. At the later point of a two-point order, a
predecessor satisfying either letter makes the corresponding
strict-past formula true and the conjunction false. This verifies
the examples and illustrates why the finite disjunction of
existential alternatives cannot be discarded.

The source's concluding use of Lemma~10 is thereby supplied in both
directions: an existential past envelope becomes a finite disjunction
of elementary past closures, and an existential future envelope
uses same-word future presentations and elementary future closures.
Universal primitive formulas and arbitrary nested regional formulas
were reduced to those cases above. Together with the full ten-lemma
interval development, this completes the stated Theorem~III.2
\cite[ch.~III, pp.~60--62, 93--94]{Kamp1968}.
\end{proof}
\end{annexitem}
